\section{Common Crawl}

Common Crawl is the largest freely available public archive of web crawl data and is released in recurring crawl snapshots \citep{baackCriticalAnalysisLargest2024,CommonCrawlGet}. It is widely used as a web-scale source of language data for search, corpus construction, NLP research, and large language model pretraining \citep{baackCriticalAnalysisLargest2024,wenzekCCNetExtractingHigh2019}. For this project, its central value lies in breadth: Common Crawl provides repeated cross-sections of general web discourse rather than data from a single platform, newspaper archive, or curated clinical corpus. However, it should not be treated as a representative survey of the web or of public opinion. The resulting data reflect what Common Crawl crawled, retained, and made available.

\begin{figure}[htbp]
	\centering
	% Source: ../reports/figures/commoncrawl_pipeline/commoncrawl_collection_pipeline.tex
	\includegraphics[width=\textwidth]{commoncrawl_pipeline/commoncrawl_collection_pipeline.pdf}
	\caption{Common Crawl collection pipeline used to construct the trend and corpus materials.}
	\label{fig:commoncrawl-pipeline}
\end{figure}

To collect the data for this study, we built a Common Crawl collection pipeline designed to extract high-quality general discourse for the analysis of diachronic lexical-semantic change in specific target terms, here ADHD and autism, against matched baseline terms. The collection workflow is illustrated in Figure~\ref{fig:commoncrawl-pipeline}. Target and baseline terms are processed together so that yearly denominators, sampling logic, and quality filters remain comparable across term groups. The pipeline uses Common Crawl's two main file formats sequentially. WET files provide compact extracted plaintext and are therefore used for large-scale term scanning and yearly prevalence denominators, but they lack the HTML structure and metadata needed for stronger validation. Candidate documents are therefore resolved to their corresponding WARC records, which preserve the archived web response, including HTML and headers. Although WARC processing is slower, it enables main-text extraction, term-survival validation, and metadata recovery. This WET-first, WARC-second design keeps the pipeline economical by reserving expensive WARC processing for candidate documents only.

The design consists of two linked tracks. The trend track uses fixed-effort annual samples to estimate how frequently target and baseline terms appear over time. The corpus track builds a larger, quality-gated document corpus for downstream NLP analysis. To reduce the risk that a small number of large websites dominate the corpus, domain caps are applied at 50 WET-validated hit rows per registered domain per Common Crawl crawl. Intermediate summaries and manifests are retained so that each year, crawl, track, and batch remains auditable. The pipeline is designed to run end-to-end on AWS EC2, using S3 as the durable storage and transfer layer for intermediate and final collection artefacts. Yearly crawl selection is deterministic: one Common Crawl snapshot is selected per year near a fixed annual anchor date and then frozen in a crawl map.

Several software choices are methodologically consequential because they affect corpus membership. WET and WARC records are parsed with \texttt{warcio}, WARC pointers are resolved through a local \texttt{pywb}-based index server, archived HTML is converted to main text with Trafilatura and Resiliparse, and post-extraction filtering uses DataTrove quality filters followed by English-language filtering with \texttt{py3langid}.

The data collection spans 13 annual Common Crawl snapshots from 2014 to 2026. In the trend track, the pipeline scanned 27.7 million WET records [provisional] and retained 78,899 WARC-validated term hits [provisional]. In the corpus track, it scanned 28.1 million WET records [provisional] and retained 42,975 analysis-ready documents [provisional]. Target-term coverage comprises 10,998 target documents [provisional], including 3,907 ADHD documents [provisional] and 8,558 autism documents [provisional]. The collection was run on an AWS \texttt{m7i-flex.large} instance, featuring 2 vCPUs, 8 GiB of RAM, and up to 12.5 Gbps network bandwidth \citep{CloudComputingServices}. On this instance type, corpus throughput was approximately one hour per million scanned WET records [provisional].


\subsection{Preprocessing}

[placeholder; leave blank for now]

\section{Target Terms}

Two target concepts were selected for diachronic lexical semantic change analysis: \textit{ADHD} and \textit{autism}. Target documents were retrieved using the matching expressions shown in Table~\ref{tab:target-patterns}. The abbreviation \texttt{ASD} was retained only when \textit{autism} occurred within $\pm$200 characters, reducing false positives from unrelated acronym use.

\begin{table}[htbp]
  \centering
  \small
  \caption{Target-term matching expressions.}
  \label{tab:target-patterns}
  \begin{tabular}{@{}ll@{}}
    \toprule
    Concept & Matching expressions \\
    \midrule
    ADHD & \texttt{\string\badhd\string\b}; \texttt{attention[-\string\s]?deficit} \\
    Autism & \texttt{\string\bautism\string\b}; \texttt{\string\bautistic\string\b}; \texttt{autism[-\string\s]?spectrum}; \texttt{\string\bASD\string\b} \\
    \bottomrule
  \end{tabular}
\end{table}

For comparison, I selected three negative, non-clinical emotion terms with sufficient coverage and interpretable usage: \textit{frustration}, \textit{sadness}, and \textit{loneliness}. These terms are not exact semantic controls for ADHD and autism; they provide a baseline for separating target-specific change from broader shifts in negative affective language in the corpus.

\section{NRC--VAD Lexicon}

The affective analyses use the NRC Valence, Arousal, and Dominance (VAD) Lexicon v2.1 \citep{mohammadNRCVADLexicon2025}. The lexicon provides real-valued ratings for approximately 55,000 English words and phrases. Scores range from 0 to 1, where higher values indicate greater valence, arousal, or dominance.

This study uses valence to estimate whether target-term contexts become more positive or negative over time, and arousal to estimate changes in emotional intensity. Dominance is present in the source lexicon but is not used. The ratings were produced through Best--Worst Scaling, where annotators are given four items (4-tuple) and asked which item is the Best (highest in terms of the property of interest) and which is the Worst (least in terms of the property of interest).
